\documentclass[12pt,a4paper]{article}
\usepackage{ugmath}
\usepackage{float}
\usepackage{placeins}
\usepackage{array}
\usepackage[labelfont=bf,labelsep=space]{caption}
\newcommand{\paren}[1]{\left( #1 \right)}
\newcommand{\alumno}{Ricardo León Martínez}
\newcommand{\materia}{Algebra Lineal I}
\newcommand{\profesor}{Claudia Reynoso Alcántara}
\newcommand{\tarea}{Examen 1}
\newcommand{\fecha}{20/3/2026}

\begin{document}

    \begin{center}
        {\large\textbf{UNIVERSIDAD DE GUANAJUATO}}\\[0.3cm]
        {\normalsize\textbf{DIVISIÓN DE CIENCIAS NATURALES Y EXACTAS}}\\
        {\normalsize\textbf{CAMPUS GUANAJUATO}}\\[1cm]

        {\Large\textbf{\tarea\ (\materia)}}\\[1cm]
    \end{center}

    \noindent
    \textbf{Nombre:} \alumno 
    \hfill 
    \textbf{Fecha:} \fecha 
    \hfill 
    \textbf{Calificación:} \rule{3cm}{0.4pt}

    \vspace{0.3cm}

    La calificación maxima del examen es de 10. Escribe lo mas claro y conciso posible.
    Usa solo la teoria vista en clase. El examen tiene una duración de 3 horas.

    \begin{mdframed}[style=mdbluebox,frametitle={Ejercicio 1}]
        Considera el siguiente sistema de ecuaciones lineales de 4
        ecuaciones lineales en 4 variables en $\mathbb{R}$.
        \begin{equation*}
            \begin{matrix}
                x_{1}+2x_{2}-x_{3}+x_{4}=3\\
                x_{1}+5x_{2}-3x_{3}=2\\
                -x_{1}-3x_{2}+2x_{3}-x_{4}=0\\
                2x_{1}+7x_{2}-4x_{3}+x_{4}=5
            \end{matrix}
        \end{equation*}
        \begin{enumerate}
            \item[(a)] Escribe la matriz aumentada del sistema.
            \item[(b)] Escribe la matriz equivalente, reducida por filas y escalonada.
            \item[(c)] Determina cuales son las variables libres y cuales son las variables no libres.
            \item[(d)] Encuentra el conjunto solucion del sistema.  
        \end{enumerate}
    \end{mdframed}

    \begin{mdframed}[style=mdbluebox,frametitle={Ejercicio 2}]
        Encuentra las ternas $(b_{1},b_{2},b_{3})\in\mathbb{R}^{3}$ para las que
        el sistema no tiene ninguna solución (justifica tu respuesta).
        \begin{equation*}
            \begin{pmatrix}
                1 & 0 & 0\\
                0 & 1 & 2\\
                1 & 1 & 2
            \end{pmatrix}
            \begin{pmatrix}
                x_{1}\\
                x_{2}\\
                x_{3}
            \end{pmatrix}
            =
            \begin{pmatrix}
                b_{1}\\
                b_{2}\\
                b_{3}
            \end{pmatrix}
        \end{equation*}
    \end{mdframed}

    \begin{mdframed}[style=mdbluebox,frametitle={Ejercicio 3}]
        Determina los valores $a\in\mathbb{R}$ para los que la siguiente
        matriz es invertible (justifica tu respuesta)
        \begin{equation*}
            \begin{pmatrix}
                1 & a & 1\\
                0 & 1 & a\\
                1 & a+1 & 2
            \end{pmatrix}
        \end{equation*}
    \end{mdframed}

    \begin{mdframed}[style=mdbluebox,frametitle={Ejercicio 4}]
        Demuestra que para todo $a\in\mathbb{R}$, $W_{a}=\{(x_{1},x_{2},x_{3},x_{4},x_{5})\in\mathbb{R}^{5}:x_{1}+x_{2}+x_{3}=ax_{5},x_{4}=0\}$
        es un subespacio vectorial de $\mathbb{R}^{5}$. Demuestra que todo
        vector en $W_{a}$ es combinación lineal de los vectores $(a,0,0,0,1)$,
        $(-1,1,0,0,0)$, $(-1,0,1,0,0)$. (justifica tu respuesta).
    \end{mdframed}

    \begin{mdframed}[style=mdbluebox,frametitle={Ejercicio 5}]
        Demuestra que $W=\{(x_{1},x_{2},x_{3},x_{4})\in\mathbb{R}^{4}:x_{1}x_{2}=0\}$
        no es subespacio vectorial de $\mathbb{R}^{4}$
    \end{mdframed}
\end{document}